\section{Conclusioni}
\label{sec:conclusioni}

Il progetto realizza i quattro compiti nella versione per una persona. Il
generatore produce griglie casuali riproducibili e le salva in JSON. Il
calcolo della chiusura determina contesto, complemento e frontiera in una
sola passata di costo lineare nel numero di celle ed è disponibile anche da
riga di comando. \cmin{} restituisce la lunghezza del cammino minimo e la
sequenza dei landmark, oppure lunghezza infinita e sequenza vuota. La
correttezza è sostenuta dalla verifica di simmetria su tutte le esecuzioni,
dal confronto con Dijkstra di tutte le varianti su griglie casuali e su
griglie costruite per forzare ricorsioni profonde, e dall'enumerazione
completa delle griglie $3\times3$ per la chiusura.

La sperimentazione mostra che il tempo di \cmin{} è spiegato quasi del tutto
dal lavoro stimato, cioè dalle chiusure calcolate moltiplicate per le celle
della griglia. Il numero di chiusure dipende dalla densità e dalla
disposizione degli ostacoli più che dalla dimensione: senza ostacoli
l'algoritmo risolve griglie $100\times100$ in poche decine di millisecondi.
Fra le ottimizzazioni, l'ordinamento della frontiera è indispensabile e la
cache conta sulle istanze difficili. Il limite propagato riduce le chiamate
del piano da \num{152548} a \num{3988} e porta le coppie completate entro
10 s da 46 a 73 su 108, senza cambiare alcuna lunghezza; la crescita resta
però combinatoria, e a $100\times100$ dal 10\% di ostacoli in su nessuna
coppia termina entro il limite di tempo.

\subsection{Limitazioni}

\begin{itemize}
  \item \textbf{Crescita combinatoria.} Il numero di chiamate non ha un
        limite polinomiale. Con 10 s e 1 GB per verso, la variante
        \code{base} è affidabile solo fino a $20\times20$ per tutte le
        densità provate e fino a $40\times40$ con il 5\% di ostacoli; il
        limite propagato estende la zona affidabile a $30\times30$ per tutte
        le densità (\cref{sec:limiti-scala}). Il piano principale resta
        dentro questa zona e non va letto come rappresentativo di griglie
        grandi e dense.
  \item \textbf{Costo lineare per chiusura.} Il calcolo della chiusura
        scandisce sempre l'intera griglia, anche quando la chiusura contiene
        poche celle.
  \item \textbf{Memoria della cache.} La cache conserva un elemento per ogni
        sottoproblema risolto, anche quando i sottoproblemi identici sono
        rari: nel caso più costoso della variante \code{base} le chiamate
        servite dalla cache sono il 13\% del totale.
  \item \textbf{Variante senza ordinamento.} Nel piano è eseguita solo sulle
        griglie più facili, perché sulle altre richiede da secondi a decine
        di secondi per esecuzione; il suo costo sulle istanze difficili è
        quindi documentato solo in modo qualitativo.
  \item \textbf{Misure dei tempi.} Il piano usa un passaggio di
        riscaldamento, un riscaldamento per esecuzione e cinque ripetizioni,
        senza un framework come JMH; il CSV conserva minimo, mediana e
        massimo ma non le misure grezze. Le tre coppie di una stessa griglia
        non sono osservazioni indipendenti, quindi correlazioni e regressioni
        descrivono il piano e non permettono inferenze generali.
  \item \textbf{Stima della memoria.} Il campo \code{memEst} conta solo le
        matrici di chiusura; il picco di heap della serie sui limiti di scala
        è un limite superiore, non una misura esatta.
  \item \textbf{Ambiente.} Il codice è compilato per Java 17 ma è stato
        eseguito solo su JDK 21.
\end{itemize}

\subsection{Sviluppi futuri}

\begin{itemize}
  \item \textbf{Chiusura incrementale.} Limitare il calcolo alla regione
        effettivamente raggiunta, invece di scandire tutta la griglia,
        renderebbe il costo di una chiamata proporzionale alla chiusura
        invece che alla griglia; per la \cref{eq:tempo} è il fattore $n$ del
        lavoro stimato.
  \item \textbf{Cache con politica di rimozione.} Limitare la cache ai
        sottoproblemi riusati più spesso ridurrebbe la memoria sulle istanze
        con pochi riusi.
  \item \textbf{Uscita anticipata dal ciclo.} Con la frontiera ordinata per
        limite inferiore, la prima cella che viola la condizione della riga
        17 implica che la violano anche le successive: il ciclo potrebbe
        terminare lì. I rami potati non ricorrono già, quindi il guadagno
        sarebbe limitato al costo del ciclo.
  \item \textbf{Misure più robuste.} Benchmark con JMH, più ripetizioni,
        salvataggio delle misure grezze ed esecuzione anche su JDK 17.
  \item \textbf{Ricostruzione del percorso.} La sequenza delle mosse si
        ottiene percorrendo il cammino libero di ogni tratto fra landmark;
        il requisito appartiene alla versione per tre persone.
\end{itemize}
